%! Author = lili
%! Date = 5/26/26

\section*{A Tutorial on Preprocessing, Estimating, and Analyzing Semantic
Networks (Christensen and Kenett 2023)}

\section{Motivation}\label{sec:motivation}

\subsection*{Why Semantic Network Analysis?}
Network science has become a popular approach to study language, memory, personality, and psychopathology~\cite{SiewEtAl2019}.
\Gls{semantic-memory} can be modelled as a network of
nodes (concepts) and edges (associations).

\Gls{verbal-fluency} is one of the most widely
used neuropsychological measures for
capturing memory retrieval~\cite{Ardila2006}.

The Core Problem (according to Christensen and
Kenett 2023) is that no standardised pipeline exists
for preprocessing \gls{verbal-fluency}
data~\cite{ChristensenKenett2023}.
Different labs make different
decisions which (can) lead to reliability threats
that compound in later analysis
steps.

Also, novice researchers lack accessible
resources to get started.

\subsection*{What Can Semantic Networks Reveal?}

The structure of \gls{semantic-memory} -- not just its content -- is a meaningful,
measurable, and clinically relevant cognitive variable.

\begin{table}[H]
    \centering
    \begin{tabular}{l p{10cm}}
    \textbf{Application Domain} & \textbf{Finding} \\
    Creative cognition & Lower average shortest path length (ASPL) linked to greater
    creative ability~\cite{KenettAnakiFaust2014} \\
        Cognition changes & Older adults show smaller clustering coefficient (CC), possibly explaining cognitive slowing~\cite{WulffHillsMata2018} \\
        Diagnostics of neurodivergence & (Asperger) autists  show ``hyper-modular'' semantic networks; rigid,
        over-partitioned knowledge~\cite{KenettGoldFaust2016} \\
    Personality (traits) & High openness to experience $\rightarrow$ more interconnected, flexi-
    ble semantic networks~\cite{ChristensenEtAl2018} \\
    Psychopathology / Diagnostics & Network structure of \gls{semantic-memory} linked to schizotypy
    and other clinical traits~\cite{ChristensenKenett2023}
    \end{tabular}
    \caption{What \glspl{semantic-network} can reveal.}\label{tab:semantic-network-reveal}
\end{table}

\section{Key Definitions}\label{sec:key-definitions}
\defin{semantic-memory}
\defin{network-graph}
\defin{verbal-fluency-task}
% TODO definitionen erfragen
We remember the definitions of:
\begin{center}
    \gls{kg}: \glsdesc{kg}
\end{center}
\begin{center}
    \gls{semantic-network}: \glsdesc{semantic-network}
\end{center}

\begin{note}
    \textbf{Node Types Depend on Task} \\
    The type of node in a \gls{semantic-network} depends on what cognitive task is being analyzed.
    For example, the node type \gls{category-exemplars} is used in \glspl{verbal-fluency-task} and the node
    type \gls{association} is used for free association tasks.
\end{note}

\subsection*{Network Science}
\defin{network-science}
According to Christensen and Kenett 2023, Network science methodologies provide a powerful computational approach for
modeling cognitive structures such as \gls{semantic-memory} and \gls{mental-lexicon}.

In the context of cognitive systems, \gls{semantic-memory} is for concept organization and the \gls{mental-lexicon}
is for word storage and word retrieval.

It is possible to measure certain aspects of \glspl{semantic-network} such as local clustering, global integration
and community structure.
Additional information about the terms (nodes) of the network can be extracted through this measuring. Measuring
methods are:
\defin{cc}
\defin{aspl}
\defin{q-modularity}


\subsection*{Verbal fluency}
\begin{center}
    ``One popular method to estimate semantic networks is via verbal fluency data.
    Verbal fluency asks participants to generate category members in 60 seconds.
    Categories can be semantic (e.g., animals) or phonological (e.g., words that start with ‘s’).''
\end{center}
(Christensen and Kenett 2023)

\frage{why-vf}

\begin{exa} To Definition
\end{exa}


\section{Problem Statement}\label{sec:problem-statement}


\section{Proposed Solution}\label{sec:proposed-solution}


\section{Network Measures}\label{sec:network-measures}


\section{Results}\label{sec:results}


\section{Limitations \& Discussion}\label{sec:limitations-&-discussion}


\section{Discussion \& Conclusion}\label{sec:discussion-&-conclusion}